\scanpage{189}\mspage{45}
% archive sheet number: 189

est plus naturel de prendre $M=e$ (ensemble ponctuel) et d'interpréter la flèche (3.2.3) comme

\[
\mathrm{Hom}(e_I,F)\longrightarrow\mathrm{Hom}(\underbrace{e_{I'}}_{\deleted{\unsure{*}}},u^*F)\tag*{(3.2.4)}
\]

et par adjonction \inserted{de $u_!,u^*$} le deuxième membre s'interprète comme $\mathrm{Hom}(u_!(e_{I'}),F)$, et la flèche est déduite de

\[
u_!e_{I'}=u_!u^*(e_I)\xrightarrow{\mathrm{adj.}}e_I \tag*{(\unsure{3.2.5})}
\]

en prenant le $\mathrm{Hom}$ dans $F$. Donc dire que (3.2.4) est iso, pour tt $F$, signifie que 3.2.5 est iso. \deleted{\unsure{Or les}} Passant aux valeurs des deux \unsure{mbs.|mbr.} en $i$, on trouve $\pi_0({}_{i\backslash}I')\overset{\sim}{\longrightarrow}1$, i.e. que $\,{}_{i\backslash}I'$ est $0$-connexe, \unsure{qed}.

% slanted
\margin{\uline{NB} Pour un énoncé \uline{complémentaire} de cofinalité cf \uline{XVIII} \unsure{prop.|par.} \unsure{4.9|4.}, 6.3 \unsure{2.} (p. 224--225, \unsure{238|228}--242) Il faudrait un énoncé qui coiffe les deux.}

\uline{Définition} 3.3. a) \deleted{\unsure{Pour que}} Quand les conditions équivalentes \inserted{de 3.2.} sont satisfaites, \inserted{\unsure{donc|i.e.} quand les $\,{}_{i\backslash}I'$ $(i\in\mathrm{Ob}\,I)$ sont $0$-conn.,} on dit que $u:\unsure{I''|I'}\longrightarrow I$ est \uline{cofinal}. On dit que $u$ est \unsure{cocofinal|cofinal} si $u^\circ:I'{}^\circ\longrightarrow I^\circ$ est cofinal, i.e. les localisées $I'_{/i}\unsure{/|,}$ de $I'$ sur $I$ $(i\in\mathrm{Ob}\,I)$ sont $0$-connexes.

b) Si $I'$ est une \uline{sous-catégorie} de $I$, on dit que $I'$ est \deleted{\unsure{cofinale dans}} $I$, si le foncteur d'inclusion est cofinal, \unsure{cocofinale} \deleted{\unsure{si le foncteur d'inclusion}} est \unsure{cocofinal|co-cofinal}.

\uline{Corollaire} 3.4. \textcircled{a} Pour que $u$ soit cofinal, \deleted{\unsure{donc}} il \uline{faut que} les cat. colocalisées $\,{}_{i\backslash}I'$ soient non vides, i.e. que l'on ait
